\section{Conclusion}
\label{sec:conclusion}

Systems for preemptible training have treated Adam's state as indivisible and the cost of recovering from a move as a
constant. Measured against a bit-identical control and a seed-to-seed noise floor, the state divides cleanly. The first
moment is a ten-step memory and can be left behind. The second moment is needed only as a per-tensor scale, which its
row and column sums supply in a hundredth of a percent of its bytes. The step counter belongs with the moments: kept
with empty moments it does harm. A first moment cannot be rebuilt at the destination, because it records the path and
not the place. Each of these statements is a prediction of a two-line model of the update and each is borne out by the
measurements.

For a scheduler the consequence is that the changeover delay is a choice. A third or a sixth of the bytes buys a
transient that cannot be told from a change of seed, and that turns a notice window which the full state misses into
one that is met several times over. The measurements here stop at \largeParams{}M parameters and one machine. The
method does not: a deterministic control, a seed band and a censored recovery time are cheap to reproduce on any
cluster, and the open question, whether the same curve holds at the sizes where the bytes hurt, is one experiment for
anyone who has the hardware.
